% 随机建模方法（与 expanda/baseline_uncertainty.py 一致）
% 本文件可用 XeLaTeX + ctex 编译，也可在定稿时翻译为英文。

\section{随机建模框架}\label{sec:stochastic}

确定性模型假设弧段运输时间和边境处理时间已知。本文在保持路径分配、流守恒、
预订运力和模态转换等一阶段决策不变的基础上，引入运输时间和边境延误两类不确定性。
采用有限场景样本近似真实分布，并以经验分位数定义成本、排放和最大完工时刻的
机会约束目标值。有限样本得到的是原随机问题的确定性近似，而不是原概率模型的精确
确定性等价形式。

\subsection{符号与路径分配}

令 $\mathcal K$ 为批次集合，$\mathcal P_k$ 为批次 $k$ 的候选路径集合，$q_k$
为货量（TEU）。$x_{kp}\in[0,1]$ 表示批次 $k$ 分配到路径 $p$ 的比例，并满足
\begin{equation}
  \sum_{p\in\mathcal P_k}x_{kp}=1,\qquad \forall k.
\end{equation}
令 $z_{kp}=\mathbf 1\{x_{kp}>0\}$ 表示路径是否实际使用。所有最大到达时刻、迟到量
和服务约束只对 $z_{kp}=1$ 的路径计算，未选择的候选路径不参与批次完成时刻。

\subsection{随机运输时间}

对弧 $a=(i,j,m)$，标称运输时间记为 $\bar T_a$，并统一使用
\begin{equation}
  \bar T_a=\frac{d_a}{v_m}.
\end{equation}
基线速度为 $(v_{\mathrm{road}},v_{\mathrm{rail}},v_{\mathrm{water}})
=(40,50,28)$ km/h。时刻表中的班次频率、首班出发时刻和发车间隔仅用于计算下一班
等待时间，不替代 $d_a/v_m$。含义或单位不明确的 \texttt{time} 字段只作为兼容
元数据读取，不直接作为逐弧运行小时数。

场景 $s$ 中的弧段时间为
\begin{equation}\label{eq:travel_random_cn}
  \tilde T_a^{(s)}=\bar T_a\xi_a^{(s)}.
\end{equation}
采用对数正态乘法扰动作为基线实验假设，其理由是运输时间非负，且延误通常形成右尾。
道路交通研究在部分稳定交通状态下支持对数正态表示，但该证据不能直接推广到铁路、
水运或所有交通状态。因此，分布类型及其参数均需进行敏感性分析
\citep{Chalumuri2014,Rakha2010,Zhao2018,Hrusovsky2018}。

令潜在乘数 $Z_m$ 满足
\begin{equation}
 \ln Z_m\sim\mathcal N\!\left(\mu_m,\sigma_m^2\right),
 \qquad \sigma_m=\sqrt{\ln(1+\mathrm{CV}_m^2)},
\end{equation}
并定义
\begin{equation}
  \xi_a^{(s)}=\min(Z_m^{(s)},c_m).
\end{equation}
代码数值求解位置参数 $\mu_m$，使
$\mathbb E[\min(Z_m,c_m)]=1$。因此 $\mathbb E[\xi_a]=1$，同时最终上界严格保持为
$c_m$。这里的 $\mathrm{CV}_m$ 决定潜在未截断分布的 $\sigma_m$；截断后实现的 CV 会略有
变化。若要求截断后的 CV 也严格等于目标值，应联合数值校准 $\mu_m$ 和 $\sigma_m$。

基线实验参数为
\begin{equation}
 (\mathrm{CV}_{\mathrm{road}},\mathrm{CV}_{\mathrm{rail}},
 \mathrm{CV}_{\mathrm{water}})=(0.15,0.10,0.20),
\end{equation}
\begin{equation}
 (c_{\mathrm{road}},c_{\mathrm{rail}},c_{\mathrm{water}})=(2.0,1.8,2.5).
\end{equation}
这些数值是基线实验假设，不是从某一篇论文直接估计的经验参数。

\subsection{随机边境延误}

令 $b=(u,v,m)$ 表示一个双边过境事件，其中 $u$ 是出境侧口岸，$v$ 是入境侧口岸。
CAREC CPMM 数据按国家、口岸、运输方式和方向报告年度平均处理时间。先将两侧均值
合并为一次过境事件的均值
\begin{equation}\label{eq:border_two_sides_cn}
  \bar D_b=\bar D_{u,\mathrm{out}}^m+
  \bar D_{v,\mathrm{in}}^m,
\end{equation}
再对完整事件只抽取一个随机延误 $\tilde D_b^{(s)}$，满足
$\mathbb E[\tilde D_b]=\bar D_b$。因此两侧不会被独立抽样或重复计时；同一场景中经过
同一方向性口岸事件的所有批次共享同一个实现值。旧网络中省略入境节点的直连弧映射
到对应的完整口岸事件。

完整事件的随机延误使用与式~\eqref{eq:travel_random_cn} 相同的截断并均值校正的
对数正态乘数。基线 CV 设为
\begin{equation}
 (\mathrm{CV}_{b}^{\mathrm{road}},\mathrm{CV}_{b}^{\mathrm{rail}},
 \mathrm{CV}_{b}^{\mathrm{water}})=(0.45,0.60,0.30),
\end{equation}
截断因子为 $c_b=3.5$。上述 CV 和截断因子都是实验设定。CAREC 年度均值只能支持
不同年份、口岸和方向存在明显差异，不能直接识别单次运输层面的 CV。CAREC 数据也
不能直接支持海港延误 CV，海运参数应由港口或班轮数据校准。对 CV=0.60 的均值为一
的对数正态变量，未校正截断前超过 3.5 倍均值的概率约为 0.56\%，而不是 0.3\%。

\subsection{固定场景集与不确定性传播}

在优化开始前一次性生成包含全部随机变量的场景集
\begin{equation}\label{eq:frozen_scenarios_cn}
  \Omega_S=\{\omega^{(1)},\ldots,\omega^{(S)}\},
\end{equation}
\begin{equation}
 \omega^{(s)}=
 \left\{\xi_a^{(s)}\ \forall a,\quad
 \tilde D_b^{(s)}\ \forall b\right\}.
\end{equation}
同一场景、同一弧或同一方向性口岸事件的实现值由所有批次和所有候选方案共享。场景集生成后冻结，
适应度评价函数内部不再调用随机数生成器。

对批次 $k$、实际使用的路径 $p$ 和场景 $s$，时钟从最早出发时刻开始：
\begin{equation}
 t_{k,0}^{(s)}=\mathrm{ET}_k.
\end{equation}
在逐弧模拟中依次加入确定性换装时间、随机边境处理时间和班次等待时间。对铁路和水运，
令 $\operatorname{dep}_a(t)$ 表示不早于 $t$ 的下一班发车时刻；对公路定义
$\operatorname{dep}_a(t)=t$。递推式为
\begin{align}
 t &\leftarrow t+\tau_j^{m_{\mathrm{in}}\to m_{\mathrm{out}}},\\
 t &\leftarrow t+\tilde D_b^{(s)},\\
 w_{k,a}^{(s)}&=\operatorname{dep}_a(t)-t,\\
 t &\leftarrow \operatorname{dep}_a(t)+\bar T_a\xi_a^{(s)}.
\end{align}
路径绝对到达时刻和运输持续时间分别为
\begin{equation}
 a_{kp}^{(s)}=t_{k,L}^{(s)},\qquad
 \tau_{kp}^{(s)}=a_{kp}^{(s)}-\mathrm{ET}_k.
\end{equation}
批次完成时刻只对实际使用的路径取最大值：
\begin{equation}
 a_k^{(s)}=\max_{p:z_{kp}=1}a_{kp}^{(s)}.
\end{equation}

基线代码将路径运力视为预订的一阶段约束，在标称计划上检查容量。随机延误影响到达、
等待、迟到与目标函数，但暂不引入随机容量。若未来把错过班次后的货物重新分配到具体
后续班次，则必须按场景重新检查班次容量，并把该规则明确为补救决策。

\subsection{场景目标函数}

场景 $s$ 的成本为
\begin{align}
 C^{(s)}={}&C^{\mathrm{transport}}+C^{\mathrm{carbon}}
 +C^{\mathrm{transship}}
 +C^{\mathrm{wait},(s)}+C^{\mathrm{border},(s)}
 +C^{\mathrm{late},(s)},\\
 C^{\mathrm{late},(s)}={}&
 \sum_k\sum_{p:z_{kp}=1}
 p_k q_kx_{kp}\,[a_{kp}^{(s)}-\mathrm{LT}_k]^+.
\end{align}
迟到罚款进入随机成本目标，但迟到本身并不自动使方案不可行。

场景排放为
\begin{equation}
 E^{(s)}=E^{\mathrm{linehaul}}+
 \sum_k\sum_a e_a^{\mathrm{wait}}q_kx_{kp}w_{k,a}^{(s)}.
\end{equation}
当基线设定 $e_a^{\mathrm{wait}}=0$ 时，排放在场景间相同。这是系统边界假设，此时排放
目标退化为确定性目标，排放置信水平不影响结果。

\subsection{迟到费率、载重和参数来源}

基线迟到费率采用 6.25 USD/TEU/hour，即 150 USD/TEU/day 除以 24；该值是按行业
滞箱费区间选取的实验基准，不再归因于 Zhao 等。当前原始工作簿中的
30--100 USD/TEU/hour 比该基准高约 4.8--16 倍，因此代码默认不使用这些值。只有在
这些逐批次费率有合同、报价或访谈记录时，才通过
\texttt{--use-input-late-penalties} 启用，并在论文中说明其来源。迟到费率只进入成本；
即使支付罚款，式~\eqref{eq:max_late_cn} 的最大迟到上限仍不可违反。

输入排放因子以 gCO$_2$/t-km 给出时，需要乘以每 TEU 平均载重。代码将该值显式设为
\texttt{--payload-tonnes-per-teu}，而不再隐藏写死。缺少本研究货物清单或运营数据时，
10 t/TEU 只能标注为实验假设，不能写成经验事实；至少对 10、14 和 15.5 t/TEU 做
敏感性分析。若使用已经按 10 t/TEU 换算的排放表，代码会按所选载重同比例缩放。

参数来源按以下优先级记录：研究网络对应的政府或运营方原始数据；官方统计或报告；
同行评议论文；合同/访谈数据；最后才是明确标注的实验假设。每次运行的场景种子、CV、
截断值、置信水平、迟到费率来源、载重和等待排放均写入
\texttt{scenario\_manifest.json}。没有来源的数值不得描述为“实测”或“采用自某文献”。

场景最大完工时刻为
\begin{equation}
 T^{(s)}=\max_{k\in\mathcal K}a_k^{(s)}.
\end{equation}
所有 $\mathrm{ET}_k$ 和 $\mathrm{LT}_k$ 必须相对于同一个时间原点计量。

\subsection{多目标机会约束模型}

一般形式为
\begin{equation}
 \min(f_C,f_E,f_T)
\end{equation}
\begin{align}
 \Pr_{\omega}\{C(x,\omega)\le f_C\}&\ge\gamma_C,\\
 \Pr_{\omega}\{E(x,\omega)\le f_E\}&\ge\gamma_E,\\
 \Pr_{\omega}\{T(x,\omega)\le f_T\}&\ge\gamma_T.
\end{align}
$f_C,f_E,f_T$ 是待最小化的可靠性上界，而不是预先指定的政策限额。若存在外部预算、
排放或时间上限，可另加 $f_C\le\bar C$、$f_E\le\bar E$、$f_T\le\bar T$。

每个批次还须满足边际准时机会约束
\begin{equation}\label{eq:ontime_marginal_cn}
 \Pr_{\omega}\{a_k(x,\omega)\le\mathrm{LT}_k\}
 \ge\gamma_{\mathrm{on}},\qquad\forall k.
\end{equation}
令
\begin{equation}
 L_k^{(s)}=[a_k^{(s)}-\mathrm{LT}_k]^+,
\end{equation}
并在全部训练场景上执行最大迟到约束
\begin{equation}\label{eq:max_late_cn}
 L_k^{(s)}\le L_k^{\max},\qquad\forall k,s.
\end{equation}
$L_k^{\max}$ 优先采用合同或 SLA。缺少业务数据时，代码采用
\begin{equation}
 L_k^{\max}=\rho(\mathrm{LT}_k-\mathrm{ET}_k),\qquad \rho=0.50,
\end{equation}
并对 $\rho\in\{0.25,0.50,0.75\}$ 进行敏感性分析。式~\eqref{eq:max_late_cn}
只保证有限训练场景满足，不等于整个不确定集合的鲁棒保证。

式~\eqref{eq:ontime_marginal_cn} 是逐批次边际约束。每个批次分别达到 90\% 准时率，
不代表全部批次同时准时的概率为 90\%。如果决策问题关注整体同时可靠性，应另设联合
机会约束
\begin{equation}
 \Pr_{\omega}\{a_k(x,\omega)\le\mathrm{LT}_k,\ \forall k\}
 \ge\gamma_{\mathrm{joint}}.
\end{equation}

\subsection{从随机模型到确定的适应度值}

固定场景集把每一个随机变量变成场景表中的具体数值。对候选方案 $x$，分别执行 $S$
次普通确定性正向模拟，得到
\begin{equation}
 \{C^{(s)}(x),E^{(s)}(x),T^{(s)}(x)\}_{s=1}^{S}.
\end{equation}
将每组结果从小到大排序。对任一目标 $Y$ 和置信水平 $\gamma$，定义经验 CCP 值
\begin{equation}\label{eq:empirical_quantile_cn}
 \hat f_Y(x)=Y^{(\lceil\gamma S\rceil)}(x).
\end{equation}
因此每个候选方案最终对应三个确定数值
\begin{equation}
 \hat F(x)=\bigl(\hat f_C(x),\hat f_E(x),\hat f_T(x)\bigr),
\end{equation}
NSGA-II 直接使用这三个数进行非支配排序，不需要在个体评价中显式建立 big-$M$ 和二元
场景违反变量。big-$M$ 形式只用于解释与机会约束文献的理论对应关系
\citep{Wang2021CCP,Li2026MOCCP}。

在固定 $\Omega_S$ 的条件下，同一个体无论评价多少次都得到完全相同的目标和约束值，
因此适应度评价中没有残余蒙特卡洛噪声。但是 NSGA-II 的初始化、交叉和变异仍具有算法
随机性；代码通过固定随机种子保证单次运行可复现，并通过多个独立算法种子评价搜索稳定性。

\subsection{样本外验证与敏感性分析}

优化完成后，应使用与训练场景独立的大型测试集评价最终 Pareto 解，不重新优化。
建议 $N_{\mathrm{test}}=5000$，报告样本外成本、排放和时间分位数、逐批次准时率、
超过 $L_k^{\max}$ 的频率、最大观察迟到以及二项比例置信区间
\citep{LuedtkeAhmed2008}。

所有 Pareto 解使用同一个测试集，以保证公平比较。算法稳定性则应通过 5--10 次完整
优化运行检验，每次改变 NSGA-II 种子和训练场景种子；仅重复测试集不能证明优化解稳定。

敏感性分析至少包括：运输时间 CV、边境延误 CV、对数正态/截断正态分布、各截断因子、
三个目标置信水平、逐批次准时置信水平、训练场景数 $S\in\{100,200,500\}$、最大迟到
比例以及非零等待排放。Pareto 超体积变化低于 2\% 只能称为经验稳定判据，比较时必须
固定归一化方法和参考点，并在多个算法种子下报告均值与标准差。

% 主要参考文献 DOI：
% Wang et al. (2021): 10.1016/j.trc.2021.103382
% Li et al. (2026): 10.1016/j.trc.2026.105699
% Zhao et al. (2018): 10.1155/2018/6051029
% Hrusovsky et al. (2018): 10.1007/s10696-016-9267-1
% Chalumuri and Asakura (2014): 10.1007/s12544-013-0111-3
% Rakha et al. (2010): 10.1080/15472450.2010.517477
% Luedtke and Ahmed (2008): 10.1137/070702928
